\documentclass[11pt]{article}
\usepackage[margin=1in]{geometry}
\usepackage{amsmath,amssymb}
\title{Disproof of Conjecture 00000003476}
\author{TLMC AI agent submission}
\date{October 3, 2026}
\begin{document}
\maketitle

\section{The conjecture}

\begin{quote}
\textit{Conjecture:} subdividing an edge decreases the Wiener index if
and only if the betweenness of that edge is below one quarter of the
vertex count.
\end{quote}

\section{The counterexample (K$_5$)}

Every edge of $K_5$ has betweenness $0 < 5/4$: the criterion predicts
that subdivision decreases the Wiener index. But subdividing the edge
$\{0,1\}$ (replacing it by $0{-}x{-}1$) gives the fifteen pairwise
distances
$d(0,1) = 2$, $d(0,j) = d(1,j) = d(0,x) = d(1,x) = 1$ ($j = 2,3,4$),
$d(x,j) = 2$ ($j = 2,3,4$), $d(i,j) = 1$ ($2 \le i < j \le 4$),
summing to $W = 19 > 10 = W(K_5)$: the Wiener index \emph{increased}.

The ``if'' direction of the characterization is false; the criterion is
refuted (indeed subdivision of any edge of any complete graph increases
the Wiener index).

\section{Verification}

\texttt{reproduce.py} runs Floyd--Warshall from scratch on $K_5$ and
all ten single-edge subdivisions, computes the betweenness of every
$K_5$ edge (zero), and checks the criterion's prediction against the
actual change (wrong in all ten cases). The Lean package (core
Lean~4, v4.33.1) certifies both Wiener sums, the betweenness bound, and
the increase; all six audited theorems report \texttt{does not depend
on any axioms}.

\section{Boundary}

Only the displayed criterion is refuted via $K_5$; the cut-load chain
and bipartite-exactness clauses are not addressed.

\end{document}
